\documentclass[11pt]{article}

\usepackage[margin=1.1in]{geometry}
\usepackage{amsmath,amssymb,amsthm}
\usepackage{mathtools}
\usepackage[round,authoryear]{natbib}
\usepackage[colorlinks=true,allcolors=blue]{hyperref}
\usepackage{microtype}

\newtheorem{claim}{Claim}
\newtheorem{proposition}{Proposition}
\theoremstyle{definition}
\newtheorem{definition}{Definition}

\DeclareMathOperator{\Var}{Var}
\DeclareMathOperator*{\argmin}{arg\,min}
\newcommand{\R}{\mathbb{R}}
\newcommand{\Sig}{\Sigma}
\newcommand{\T}{^{\mathsf{T}}}
\newcommand{\smin}{\sigma_{\min}}
\newcommand{\smax}{\sigma_{\max}}

\title{\textbf{Schur Damping for Perpetual Demand Lending Pools}}
\author{Peter Cotton}
\date{June 2026}

\begin{document}
\maketitle

\begin{abstract}
Perpetual demand lending pools (PDLPs) must decide whether to operate as one
pool or several. \citet{chitra2025} resolve this with the \emph{undamped} Schur
complement of the asset covariance and a binary spectral threshold: merge the
pools if a single condition number test passes, keep them apart otherwise. We
observe that the merged and isolated pools are the two endpoints, $\gamma=1$ and
$\gamma=0$, of the Schur-complementary damping family of \citet{cotton2024}, and
we propose occupying its interior. A \emph{partially merged} pool hedges against
$A-\gamma\,B C^{-1} B\T$, with the coupling fraction $\gamma$ set not by a
threshold but by the statistical reliability of the cross-pool covariance block
$B$. We give the closed-form $\gamma^\star$ implied by viewing the choice as a
James--Stein shrinkage problem \citep{cotton2026}. Because PDLPs live in an acute
undersampled regime---short on-chain return histories, fast rebalancing, asset
counts that rival or exceed the number of return samples---$\gamma^\star$ is
pulled toward $0$: the theory predicts pools should stay \emph{mostly} split,
merging only as cross-pool coupling proves both strong and well-sampled. The
result is a graded, data-driven target-weight mechanism that degrades gracefully
where the binary test flips discretely, and it supplies precisely the pool
covariance that GMX~V2's dynamic-pricing upgrade is noted to omit.
\end{abstract}

\section{Introduction}

Perpetual futures are the most liquid crypto markets, and a growing share trades
on decentralized venues---GMX, Jupiter, Hyperliquid---that fund market-maker
leverage through on-chain liquidity pools. \citet{chitra2025} formalize this
mechanism as a \emph{perpetual demand lending pool} (PDLP): liquidity providers
deposit assets into a pool with a target composition, traders borrow against it
to open leveraged positions, and arbitrageurs keep the pool near target. The
authors give conditions under which a PDLP's liquidity providers can be delta
hedged, partially explaining why these pools have grown to billions of dollars in
assets.

A recurring architectural question, raised explicitly in their Appendix~B, is
when several pools should be combined into one. GMX itself migrated from a single
pool to a multi-pool system; the inverse question---when to re-aggregate---is
decided in \citet{chitra2025} by partitioning the asset covariance into blocks and
comparing a pool's own covariance against its Schur complement. The comparison is
binary: a spectral inequality either holds, and one merged pool delta hedges
better, or it fails, and the pools should stay apart.

This note makes a single observation and develops its consequence. The two
options being compared---the isolated pool and the merged pool---are not two
unrelated architectures. They are the $\gamma=0$ and $\gamma=1$ endpoints of the
damped Schur complement that interpolates hierarchical risk parity and minimum
variance \citep{cotton2024,cotton2022}. The interesting regime for a PDLP is the
interior, $0<\gamma<1$: a pool that is \emph{partially} merged, sharing risk
across pools in proportion to how reliably that shared risk can be estimated. In
the data-poor conditions typical of DeFi, that proportion is small, and the
interior solution differs materially from either endpoint.

\section{The pool-merging decision is the $\gamma$-bridge endpoints}

We recall the setup of \citet[App.~B]{chitra2025}. A PDLP holds risky assets with
covariance $\Sig \succ 0$. To ask whether two pools should be one, partition the
assets into a pool-1 block of size $m_1$ and a pool-2 block of size $m_2$ with
disjoint supports, $m_1+m_2=n$, and write
\begin{equation}
  \Sig = \begin{bmatrix} A & B \\ B\T & C \end{bmatrix},
  \qquad A\in\R^{m_1\times m_1},\ C\in\R^{m_2\times m_2},\ B\in\R^{m_1\times m_2},
\end{equation}
with $A,C$ nonsingular and positive definite. Here $A$ is pool~1's own
covariance, $C$ is pool~2's, and $B$ is the \emph{cross-pool} block. For
mean-variance delta hedging the relevant objects are the conditional covariances,
the Schur complements
\begin{equation}
  \Sig/C = A - B\,C^{-1}B\T, \qquad \Sig/A = C - B\T A^{-1} B .
\end{equation}
A merged pool incorporates information from both blocks and hedges against
$\Sig/C$ (resp.\ $\Sig/A$); an isolated pool hedges against $A$ (resp.\ $C$)
alone, as though the other pool did not exist. \citet{chitra2025} prove a
sufficient condition for one pool to be better than two.

\begin{claim}[{\citealp[Claim~B.1]{chitra2025}}]
\label{claim:b1}
If $\smax(A) < \smin(\Sig/A)$ and $\smax(B\T)<\smin(\Sig/C)$ then the single
merged pool delivers better delta-hedged returns than the two isolated pools.
\end{claim}

The authors note these conditions are ``similar to hierarchical risk-parity
methods'' \citep{cotton2024,deprado2016}. The similarity is exact, and it is the
point of departure for what follows: the isolated pool keeps only its diagonal
block $A$, while the merged pool conditions fully through $A - BC^{-1}B\T$. These
are the two ends of one dial.

\section{Damped pool merging}

Following \citet{cotton2024}, interpolate the coupling. For $\gamma\in[0,1]$
define the \emph{damped sub-covariance}
\begin{equation}
\label{eq:damped}
  A(\gamma) \;=\; A - \gamma\, B\,C^{-1}B\T,
  \qquad
  C(\gamma) \;=\; C - \gamma\, B\T A^{-1} B ,
\end{equation}
and let a \emph{partially merged} PDLP delta hedge pool~1 against $A(\gamma)$ and
pool~2 against $C(\gamma)$.

\begin{proposition}[Endpoints]
$A(0)=A$ recovers the isolated pool, and $A(1)=\Sig/C$ recovers the fully merged
pool of \citet{chitra2025}. The single-pool-versus-two-pools decision of
Claim~\ref{claim:b1} is the comparison of the two endpoints $\gamma\in\{0,1\}$.
\end{proposition}

\noindent The proposition is immediate from \eqref{eq:damped}, but it reframes the
problem. Rather than testing which endpoint is better, choose $\gamma$. And
$\gamma$ has a meaning: it is how much of the estimated cross-pool coupling to
trust.

\paragraph{Choosing $\gamma$.} The right $\gamma$ is the reliability of the
estimated coupling, a James--Stein / Wiener ratio of signal to signal-plus-noise
rather than a style preference \citep{cotton2026}. In the two-block case with
conditional correlation $\rho$ estimated from $n$ return observations,
\begin{equation}
\label{eq:gammastar}
  \gamma^\star \;=\; \frac{(n-2)\,\rho^2}{(n-2)\,\rho^2 + (1-\rho^2)} .
\end{equation}
Strong, well-sampled cross-pool coupling sends $\gamma^\star\to 1$ and the merge
deserves trust; weak or undersampled coupling sends $\gamma^\star\to 0$ and the
divide-and-conquer split was right. The estimate $\gamma^\star$ tracks the
validation-tuned optimum with no tuning, on portfolios, crypto correlation, and
simulated processes \citep{cotton2024,cotton2026}.

\section{The undersampled regime is where DeFi lives}

The binary test of Claim~\ref{claim:b1} silently assumes $B$ is known. In a PDLP
it is the least-known part of the model. The pools are disjoint by construction,
so $B$ is a cross-\emph{block} covariance---the entries with the fewest joint
observations and the most estimation noise. On-chain return histories are short,
pools rebalance on the order of hours, and the number of pooled assets can rival
or exceed the number of return samples $n$. Trusting $B$ completely (merge)
overfits noise into the hedge; discarding it (split) throws away genuine
diversification. The reliability-weighted answer is to trust $B$ in proportion to
how well it is measured, which is exactly \eqref{eq:gammastar}.

Small $n$ pulls $\gamma^\star$ toward $0$. The theory therefore predicts that
PDLP pools should remain \emph{mostly} split, merging incrementally only as the
cross-pool coupling proves both strong (large $\rho$) and stable (large $n$).
This is a graded, falsifiable prediction. The threshold test of
Claim~\ref{claim:b1} cannot make it: it flips from ``split'' to ``merge'' the
instant a spectral inequality is met and is silent on the long approach to that
boundary, which is where a young, data-poor pool actually operates.

This is the same identity that drives \citet{cotton2026}: the damping a portfolio
needs when assets outnumber returns is, term for term, the damping a spatial
pseudo-likelihood needs when stations outnumber observations. A perpetual demand
lending pool is an unusually sharp instance of the outnumbered regime, which is
why its architecture question lives in the interior of the dial rather than at its
ends.

\section{A reliability-weighted target-weight mechanism}

\citet{chitra2025} maintain a pool near a target portfolio through a
target-weight mechanism, and close Appendix~B by noting that GMX~V2's
dynamic-pricing upgrade ``does not take into account pool asset covariance.''
Equation~\eqref{eq:damped} supplies exactly that covariance, at a
reliability-weighted amount. Concretely, a damped PDLP hedges pool~1 against
$A(\gamma^\star)$ rather than against either $A$ or $\Sig/C$, and a target-weight
mechanism can set inter-pool collateral sharing proportional to $\gamma^\star$:
\begin{equation}
  \text{shared coupling} \;=\; \gamma^\star \cdot B\,C^{-1}B\T,
  \qquad
  \gamma^\star \in [0,1].
\end{equation}
When the cross-pool block is well estimated the mechanism behaves like a single
diversified pool; when it is not, it falls back continuously toward isolated
pools, degrading gracefully rather than discretely. No new primitive is required:
$A(\gamma^\star)$ is the augmented sub-covariance of \citet{cotton2024}, and the
only added ingredient is a running estimate of $\rho$ and $n$ for each pool pair,
both available on chain.

\section{Discussion}

The contribution is deliberately small. \citet{chitra2025} already reach for the
Schur complement and already note its kinship with hierarchical risk parity; what
is missing is the damping parameter that turns a binary architectural test into a
continuous, self-tuning one. The same move---occupying the interior of a dial that
a field has only used at its endpoints---recurs in portfolio construction (HRP
versus minimum variance), in spatial statistics (base-shrunk versus full Vecchia
conditionals), and in optimization (KFAC's binary block-coupling choices)
\citep{cotton2024,cotton2026}. DeFi lending pools are, if anything, the cleanest
case for damping, because their estimation problem is the most severe. We leave to
future work an empirical comparison of damped versus binary pool architectures on
historical PDLP data, and the connection between the dynamic-pricing model and the
cost of delta hedging that \citet{chitra2025} flag as open.

\begin{thebibliography}{9}

\bibitem[Chitra et al.(2025)]{chitra2025}
T.~Chitra, T.~Diamandis, N.~Sheng, L.~Sterle, and K.~Yusubov.
\newblock Perpetual Demand Lending Pools.
\newblock arXiv:2502.06028, 2025.

\bibitem[Cotton(2022)]{cotton2022}
P.~Cotton.
\newblock Schur Complementary Portfolios Fix Hierarchical Risk Parity.
\newblock \emph{Geek Culture} (Medium), 2022.

\bibitem[Cotton(2024)]{cotton2024}
P.~Cotton.
\newblock Schur Complementary Allocation: A Unification of Hierarchical Risk
Parity and Minimum Variance Portfolios.
\newblock arXiv:2411.05807, 2024.

\bibitem[Cotton(2026)]{cotton2026}
P.~Cotton.
\newblock Two Sides of Schur Damping: High-Dimensional Pseudo-Likelihoods and
Portfolio Allocation.
\newblock arXiv:2606.14798, 2026.

\bibitem[L\'opez de Prado(2016)]{deprado2016}
M.~L\'opez de Prado.
\newblock Building Diversified Portfolios that Outperform Out of Sample.
\newblock \emph{Journal of Portfolio Management}, 42(4), 2016.

\end{thebibliography}

\end{document}
